\section{Kiến trúc kỹ thuật}
\label{sec:technical-architecture}
\begingroup
\setlength{\emergencystretch}{3em}
\captionsetup[table]{name=Bảng}
\clubpenalty=10000
\widowpenalty=10000

OneSpace tích hợp video, âm nhạc, tin tức và trò chơi trong cùng một ứng dụng Android. Khác với việc xây dựng bốn giao diện độc lập, sự tích hợp này đòi hỏi các phân hệ phối hợp tài nguyên, duy trì trạng thái nội dung và phản ứng nhất quán khi người dùng chuyển hoạt động. Chẳng hạn, việc mở một bài báo trong khi nghe nhạc chỉ thay đổi ngữ cảnh giao diện; trong khi bắt đầu phát một video có tiếng lại làm phát sinh yêu cầu phân xử âm thanh.

Chương này trình bày cơ sở thiết kế kiến trúc cho các tình huống trên. Các quyết định được xem xét theo bốn tiêu chí: tính liên tục của trải nghiệm, khả năng thích ứng với vòng đời Android, mức sử dụng tài nguyên và khả năng bảo trì. Nội dung lần lượt phân tích ràng buộc môi trường di động, lựa chọn công nghệ, tổ chức các thành phần và mô hình quản lý trạng thái khi thực thi.

Kiến trúc đề xuất sử dụng mô hình ba lớp Presentation--Domain--Data, kết hợp mô hình Model--View--ViewModel (MVVM) và luồng dữ liệu một chiều (Unidirectional Data Flow, UDF). Mã nguồn được tổ chức thành các mô-đun Gradle có trách nhiệm xác định; mô-đun ứng dụng đảm nhiệm khởi tạo và kết nối các thành phần. Đây là cách áp dụng các khuyến nghị Android vào nhu cầu của OneSpace, thay vì coi một cấu trúc mô-đun cụ thể là yêu cầu bắt buộc của nền tảng~\cite{android-architecture,android-modularization}.

\textbf{Baseline kiến trúc thống nhất ngày 24/10/2026.} Phần 4.1--4.4 xác định kiến trúc đích; ma trận cuối tài liệu phân biệt phần đã có trong code và phần cần triển khai tiếp. Việc chốt kiến trúc không có nghĩa mọi chức năng đã hoàn thành. Phạm vi hiện thực hóa trong bản nguyên mẫu và các giới hạn được trình bày riêng tại mục~\ref{sec:architecture-scope}.

Phạm vi kiến trúc phục vụ CO3043: phân biệt môi trường mobile và desktop (L.O.1.1), giải thích luồng thực thi và vòng đời (L.O.1.2), thiết kế module với ranh giới rõ ràng (L.O.1.3), lựa chọn framework phù hợp (L.O.4.1). Không mở rộng sang microservices, nhiều nền tảng hoặc game engine khi chưa có nhu cầu trong bài.

\begingroup
\small
\renewcommand{\arraystretch}{1.2}
\setlength{\tabcolsep}{3pt}
\begin{longtable}{>{\RaggedRight\arraybackslash}p{0.26\linewidth}>{\RaggedRight\arraybackslash}p{0.68\linewidth}}
\caption{Technical baseline của OneSpace}
\label{tab:baseline}\\
\toprule
\textbf{Hạng mục} & \textbf{Quyết định} \\
\midrule
\endfirsthead
\toprule
\textbf{Hạng mục} & \textbf{Quyết định} \\
\midrule
\endhead
\bottomrule
\endfoot
Nền tảng triển khai & Android Native; không tuyên bố hỗ trợ iOS \\
Ngôn ngữ/giao diện & Kotlin, Jetpack Compose, Material 3 \\
Thiết bị & minSdk 26; RAM 3 GB là mục tiêu kiểm thử, chưa phải kết quả benchmark \\
Kiến trúc & Ba layer Presentation--Domain--Data, MVVM và UDF \\
Tổ chức & Gradle multi-module: app, domain, feature/*, data/*, core/* \\
Điều hướng & Một Activity, một NavController, một root NavHost; nested graphs theo feature \\
Media & Media3; tối đa một video player và một audio player dùng chung \\
Âm thanh & Một nguồn nội dung chính có tiếng tại một thời điểm; video mute có thể cùng tồn tại với nhạc \\
Lưu trữ & Room cho dữ liệu có cấu trúc/tiến độ; DataStore cho preferences và metadata phiên nhỏ \\
Chạy nền & MediaSessionService là kiến trúc đích của Music; WorkManager dành cho công việc trì hoãn \\
Trạng thái hiện tại & Catalog mock, auth demo; phát sample media thật; các phần còn thiếu nằm trong ma trận triển khai \\
\end{longtable}
\endgroup

\subsection{Ràng buộc của môi trường di động}
\label{sec:mobile-constraints}

\subsubsection{Mobile App so với Web App và Desktop App}
\label{sec:mobile-vs-desktop}

Yêu cầu L.O.1.1 của học phần yêu cầu phân tích sự khác biệt kiến trúc giữa ứng dụng di động và ứng dụng máy tính. Bảng~\ref{tab:mobile-desktop} so sánh ba loại ứng dụng theo các tiêu chí ảnh hưởng trực tiếp đến thiết kế của OneSpace.

\begin{longtable}{>{\raggedright\arraybackslash}p{2.6cm} >{\raggedright\arraybackslash}p{4.4cm} >{\raggedright\arraybackslash}p{4.2cm} >{\raggedright\arraybackslash}p{4.2cm}}
\caption{So sánh Native Mobile App, Web App và Desktop App}
\label{tab:mobile-desktop}\\
\toprule
\textbf{Tiêu chí} & \textbf{Native Mobile App} & \textbf{Web App} & \textbf{Desktop App} \\
\midrule
\endfirsthead
\toprule
\textbf{Tiêu chí} & \textbf{Native Mobile App} & \textbf{Web App} & \textbf{Desktop App} \\
\midrule
\endhead
\bottomrule
\endfoot
Hiệu năng & Tốt, truy cập API hệ điều hành trực tiếp & Phụ thuộc trình duyệt & Thường có nhiều tài nguyên hơn \\
RAM, CPU & Giới hạn theo thiết bị & Phụ thuộc trình duyệt và thiết bị & Thường ít bị giới hạn \\
Pin & Cần tối ưu nghiêm ngặt & Cũng cần tối ưu trên điện thoại & Ít bị ảnh hưởng khi cắm nguồn \\
Phát nhạc nền & MediaSession / Service & Có thể bị giới hạn bởi trình duyệt & Tương đối linh hoạt \\
Chạy game & Native rendering & Canvas / WebGL & Có thể dùng GPU mạnh \\
Offline & Room, DataStore, file cache & Service Worker, IndexedDB & Local filesystem / database \\
Vòng đời & Do hệ điều hành kiểm soát chặt & Phụ thuộc vòng đời tab/trình duyệt & Ít bị hạn chế hơn \\
Thông báo & Android Notification API & Web Notifications khi được hỗ trợ & OS notification \\
Phù hợp dự án & \textbf{Rất phù hợp} & Hạn chế ở tích hợp nền & Không phù hợp mục tiêu mobile \\
\end{longtable}

Kết luận: ứng dụng tiêu dùng nội dung giải trí đa phương tiện cần phát nhạc nền ổn định, xử lý vòng đời nghiêm ngặt và thông báo tức thì~--~đây là những khả năng chỉ native mobile app đáp ứng đủ cho phạm vi môn học.

\subsubsection{Phân hệ nào tiêu thụ nhiều tài nguyên}
\label{sec:resource-consumption}

Bảng~\ref{tab:resource-consumption} đánh giá tương đối mức tiêu thụ tài nguyên theo phân hệ (định tính, không phải benchmark), làm cơ sở cho chính sách tối ưu ở mục~\ref{sec:resource-policy}.

\begin{longtable}{>{\raggedright\arraybackslash}p{1.6cm} >{\centering\arraybackslash}p{1.3cm} >{\centering\arraybackslash}p{1.8cm} >{\centering\arraybackslash}p{1.8cm} >{\raggedright\arraybackslash}p{8.0cm}}
\caption{Mức tiêu thụ tài nguyên theo phân hệ (đánh giá định tính)}
\label{tab:resource-consumption}\\
\toprule
\textbf{Phân hệ} & \textbf{Pin} & \textbf{Dữ liệu mạng} & \textbf{CPU/RAM} & \textbf{Biện pháp} \\
\midrule
\endfirsthead
\toprule
\textbf{Phân hệ} & \textbf{Pin} & \textbf{Dữ liệu} & \textbf{CPU/RAM} & \textbf{Biện pháp} \\
\midrule
\endhead
\bottomrule
\endfoot
Video & Cao & Rất cao & Cao & Adaptive bitrate, giới hạn preloading, giải mã phần cứng \\
Music & Thấp--TB & Thấp--TB & Thấp & Audio cache, phát với màn hình tắt, tránh polling \\
News & Thấp & Thấp--TB & Thấp & Cache bài viết, lazy loading ảnh \\
Games & TB--Cao & Thấp nếu offline & Phụ thuộc game & Giảm FPS, dừng game loop khi pause \\
\end{longtable}

Video thường tiêu thụ dữ liệu di động nhiều nhất và là một trong những nguồn hao pin lớn nhất. Tuy nhiên, game có đồ họa phức tạp vẫn có thể tốn pin hơn video~--~đây là lý do feature F-GAME-15 (kiểu flappy) được chấm Effort 2 thay vì 1 tại bảng~\ref{tab:feature-pool}.

\subsubsection{Chính sách tối ưu tài nguyên}
\label{sec:resource-policy}

Nhóm xây dựng một \texttt{ResourcePolicyManager} điều phối các quy tắc xử lý theo tình huống:

\begin{longtable}{>{\raggedright\arraybackslash}p{5.2cm} >{\raggedright\arraybackslash}p{9.8cm}}
\caption{Chính sách xử lý tài nguyên theo tình huống}
\label{tab:resource-policy}\\
\toprule
\textbf{Tình huống} & \textbf{Quy tắc xử lý} \\
\midrule
\endfirsthead
\toprule
\textbf{Tình huống} & \textbf{Quy tắc xử lý} \\
\midrule
\endhead
\bottomrule
\endfoot
Video chuyển thành Floating & Chỉ giữ một video decoder đang hoạt động \\
Music chạy nền & Dùng MediaSessionService, không giữ UI chạy \\
Game chuyển tab & Pause game loop, lưu game state \\
News ngoài màn hình & Ngừng các tác vụ UI không cần thiết \\
RAM thấp & Giảm prefetch, giải phóng player và ảnh không dùng \\
Bật Data Saver & Giảm chất lượng streaming, hạn chế tải trước \\
Battery Saver & Giảm xử lý nền không thiết yếu \\
\end{longtable}

Đặc biệt, khi game không còn hiển thị, không giữ vòng lặp render 60~FPS chỉ để người dùng có thể quay lại nhanh~--~điều này gây lãng phí trên thiết bị RAM 3~GB của persona Quốc Huy (mục~\ref{sec:persona-quochuy}).

\subsubsection{Offline và mạng không ổn định}
\label{sec:offline-design}

Bảng~\ref{tab:offline-behavior} định nghĩa hành vi từng phân hệ khi mất mạng và khi có mạng trở lại:

\begin{longtable}{>{\raggedright\arraybackslash}p{1.6cm} >{\raggedright\arraybackslash}p{6.8cm} >{\raggedright\arraybackslash}p{6.8cm}}
\caption{Hành vi offline theo phân hệ}
\label{tab:offline-behavior}\\
\toprule
\textbf{Phân hệ} & \textbf{Khi mất mạng} & \textbf{Khi có mạng trở lại} \\
\midrule
\endfirsthead
\toprule
\textbf{Phân hệ} & \textbf{Khi mất mạng} & \textbf{Khi có mạng trở lại} \\
\midrule
\endhead
\bottomrule
\endfoot
Video & Phát phần còn buffer/cache; nếu không đủ thì chờ mạng & Retry hoặc resume từ vị trí trước đó \\
Music & Phát bài đã tải xuống; streaming chỉ phát được phần buffer & Tiếp tục streaming \\
News & Đọc bài đã cache, bookmark vẫn hoạt động & Cập nhật nội dung từ server \\
Games & Game offline vẫn chơi được (F-GAME-12/13) & Đồng bộ điểm số nếu có tài khoản \\
\end{longtable}

Lưu ý quan trọng: \textbf{cache không đồng nghĩa với download}. Streaming cache có thể bị xóa và không bảo đảm phát ngoại tuyến toàn bộ nội dung. Nếu muốn chức năng Downloaded Music (F-MUS-12), cần thiết kế lưu trữ và quản lý file riêng~--~tính năng này đã được xếp vào nhóm Thankless (hoãn v1) tại bảng~\ref{tab:feature-pool}.

\subsubsection{Cấu hình tối thiểu hỗ trợ}
\label{sec:min-config}

\begin{longtable}{>{\raggedright\arraybackslash}p{2.6cm} >{\raggedright\arraybackslash}p{12.4cm}}
\caption{Cấu hình tối thiểu đề xuất}
\label{tab:min-config}\\
\toprule
\textbf{Thành phần} & \textbf{Quyết định} \\
\midrule
\endhead
\bottomrule
\endfoot
Android & 8.0+ (API 26+) \\
RAM & 3~GB mục tiêu (khớp persona Quốc Huy, thiết bị phổ thông) \\
Storage & Khoảng 200~MB trống cho cài đặt và cache MVP \\
Internet & Cần cho streaming và news mới \\
Offline & News cached, downloaded music, minigame offline \\
iOS & \textbf{Không hỗ trợ trong MVP}~--~phân tích iOS có ở mục~\ref{sec:platform-selection} nhưng nhóm không code bản iOS \\
\end{longtable}

API 26 và RAM 3~GB là mục tiêu hỗ trợ đề xuất, không phải mức tối thiểu bắt buộc của Jetpack Compose. Cần kiểm thử trên emulator hoặc thiết bị tương ứng (danh sách test case ở mục~\ref{sec:lifecycle-state}).

\subsection{Lựa chọn nền tảng và công nghệ}
\label{sec:platform-selection}

\subsubsection{Phương án đã cân nhắc}
\label{sec:platform-comparison}

\begingroup
\small
\renewcommand{\arraystretch}{1.2}
\setlength{\tabcolsep}{3pt}
\begin{longtable}{>{\RaggedRight\arraybackslash}p{0.17\linewidth}>{\RaggedRight\arraybackslash}p{0.28\linewidth}>{\RaggedRight\arraybackslash}p{0.22\linewidth}>{\RaggedRight\arraybackslash}p{0.23\linewidth}}
\caption{Đối chiếu các nền tảng và framework}
\label{tab:platform-comparison}\\
\toprule
\textbf{Phương án} & \textbf{Điểm phù hợp} & \textbf{Chi phí trong bài} & \textbf{Kết luận} \\
\midrule
\endfirsthead
\toprule
\textbf{Phương án} & \textbf{Điểm phù hợp} & \textbf{Chi phí trong bài} & \textbf{Kết luận} \\
\midrule
\endhead
\bottomrule
\endfoot
Android Native/Kotlin & Trực tiếp dùng lifecycle, Media3, media session và audio focus; nhóm có kinh nghiệm & Chỉ triển khai Android & Chọn \\
iOS Native/Swift & Tích hợp hệ thống iOS trực tiếp & Thêm nền tảng và công cụ; không có bản iOS để kiểm thử & Phân tích, không triển khai \\
Flutter & Chia sẻ UI và logic đa nền tảng & Cần kiểm tra plugin và phần tích hợp Android cho từng yêu cầu media & Không chọn trong phạm vi hiện tại \\
React Native & Chia sẻ UI/logic bằng JS/TS & Phụ thuộc thư viện/native integration cho media và background & Không chọn \\
Kotlin Multiplatform & Có thể chia sẻ logic Kotlin & Vẫn cần chiến lược UI/media theo nền tảng; tăng phạm vi kiểm thử & Cân nhắc sau nếu có yêu cầu iOS \\
\end{longtable}
\endgroup

Không dùng bảng này để khẳng định một framework luôn nhanh hơn framework khác hoặc mọi cross-platform app đều cần tự viết native bridge.

\subsubsection{Technology stack và mức triển khai}
\label{sec:tech-stack}

Bảng~\ref{tab:tech-stack} liệt kê công nghệ theo vai trò trong kiến trúc và mức triển khai tại bản nguyên mẫu hiện tại.

\begingroup
\small
\renewcommand{\arraystretch}{1.2}
\setlength{\tabcolsep}{3pt}
\begin{longtable}{>{\RaggedRight\arraybackslash}p{0.23\linewidth}>{\RaggedRight\arraybackslash}p{0.32\linewidth}>{\RaggedRight\arraybackslash}p{0.37\linewidth}}
\caption{Technology stack và mức triển khai}
\label{tab:tech-stack}\\
\toprule
\textbf{Vai trò} & \textbf{Công nghệ chốt} & \textbf{Hiện trạng} \\
\midrule
\endfirsthead
\toprule
\textbf{Vai trò} & \textbf{Công nghệ chốt} & \textbf{Hiện trạng} \\
\midrule
\endhead
\bottomrule
\endfoot
UI/design system & Compose, Material 3, font đóng gói & Có \\
Navigation & Navigation Compose typed routes + Kotlin Serialization & Có \\
State/DI & ViewModel, StateFlow, lifecycle-aware collection, Hilt/KSP & Có \\
Domain & Kotlin/JVM, coroutines/Flow cho contracts & Có \\
Media & Media3 ExoPlayer, adapter thực thi MediaSessionPort & Có phát sample stream thật \\
Structured local data & Room & Có 6 bảng: video progress, watchlist, điểm, nhiệm vụ hằng ngày, âm thanh trích, nội dung tải (schema v5) \\
Preferences & DataStore & Có \\
Remote API & Retrofit + OkHttp + Kotlin Serialization & Chọn cho giai đoạn remote; chưa có client/backend thật \\
Artwork/ảnh & Coil khi tích hợp ảnh từ nguồn thật & Chưa tích hợp; hiện dùng minh họa mock \\
Music background & MediaSessionService + MediaController & Bắt buộc cho hành vi background đích; chưa triển khai \\
Deferred sync & WorkManager & Dành cho sync/retry trì hoãn; chưa có worker \\
Mini game & Logic Kotlin + Compose/Canvas phù hợp game nhỏ & Có: board game và mini-game nhịp chạy thật, tích hợp thưởng điểm \\
Verification & JUnit, Compose UI tests, verifyArchitecture & Có baseline; mỗi yêu cầu mới cần bằng chứng riêng \\
\end{longtable}
\endgroup

Không thêm dependency, empty service hoặc empty use case chỉ để bảng stack trông đầy đủ. Những mục chưa có phải giữ trạng thái dự kiến cho đến khi có code và test.

\subsection{Kiến trúc ba lớp và chia module}
\label{sec:layered-architecture}

\begin{guide}
\guidepart{Nội dung}
\begin{compactitem}
    \item Presentation Layer: UI component và quản lý trạng thái (Provider, Bloc, Redux, \ldots tùy framework chọn ở mục~\ref{sec:platform-selection}).
    \item Business Logic Layer (Domain): xác thực người dùng, kiểm tra dữ liệu form, quy tắc nghiệp vụ riêng cho Video, Audio, News và Game.
    \item Data Layer: lưu trữ cục bộ (Local Storage) và nguồn dữ liệu từ xa (Remote Data Sources).
    \item Cách chia module: mỗi phân hệ là một module riêng, cùng các module dùng chung (tài khoản, trình phát media, mạng).
\end{compactitem}
\guidepart{Câu hỏi cần trả lời}
\begin{compactitem}
    \item Các lớp phụ thuộc nhau theo chiều nào? Lớp Domain có phụ thuộc vào framework giao diện không?
    \item Dữ liệu nào được lưu cục bộ, dữ liệu nào luôn lấy từ server, và hai nguồn được đồng bộ ra sao?
    \item Muốn thêm một phân hệ mới (ví dụ podcast) thì cần thay đổi những module nào?
\end{compactitem}
\guidepart{Hình/bảng nên có}
Sơ đồ kiến trúc 3 lớp, sơ đồ phụ thuộc giữa các module.
\end{guide}

\subsection{Vòng đời ứng dụng và quản lý trạng thái}
\label{sec:lifecycle-state}

\subsubsection{Luồng thực thi (Execution Flow)}
\label{sec:execution-flow}

Luồng khởi động của OneSpace:

\begin{enumerate}
\item \textbf{Launch app}: Android khởi chạy MainActivity, OneSpaceApplication khởi tạo Hilt dependencies.
\item \textbf{Load preferences and saved session}: đọc DataStore (cài đặt người dùng) và Room/DataStore (phiên hoạt động trước đó).
\item \textbf{Kiểm tra phiên đăng nhập}: nếu có phiên hợp lệ $\rightarrow$ vào Home trực tiếp; nếu không $\rightarrow$ vào Guest Mode hoặc màn Đăng nhập.
\item \textbf{Home Screen}: hiển thị nội dung gợi ý theo sở thích; người dùng chọn phân hệ (Video, Music, News, Games).
\item \textbf{SessionCoordinator} lưu hoặc khôi phục trạng thái phân hệ tương ứng.
\end{enumerate}

Nhóm khuyên \textbf{Guest Mode} trong MVP (tương ứng F-COM-18) để người dùng trải nghiệm đủ bốn phân hệ mà không bị phụ thuộc backend. Authentication là module mở rộng nếu assignment không yêu cầu tài khoản bắt buộc.

\subsubsection{Quy tắc giữ trạng thái}
\label{sec:state-rules}

Điểm rất quan trọng: \textbf{chuyển tab trong app khác với đưa toàn bộ ứng dụng xuống background}. Chuyển \texttt{Video $\rightarrow$ Tin} vẫn nằm trong cùng Activity, nên Floating Video tiếp tục phát. Nhưng bấm Home đưa Activity xuống nền thì MVP \textbf{pause video} để tiết kiệm pin, trừ khi người dùng kích hoạt chế độ riêng.

Bảng~\ref{tab:state-rules} định nghĩa trạng thái từng loại nội dung trong ba tình huống:

\begin{longtable}{>{\raggedright\arraybackslash}p{1.6cm} >{\raggedright\arraybackslash}p{4.6cm} >{\raggedright\arraybackslash}p{5.2cm} >{\raggedright\arraybackslash}p{4.6cm}}
\caption{Quy tắc giữ trạng thái theo ba tình huống}
\label{tab:state-rules}\\
\toprule
\textbf{Nội dung} & \textbf{Chuyển tab trong app} & \textbf{App vào background} & \textbf{Process bị kill} \\
\midrule
\endfirsthead
\toprule
\textbf{Nội dung} & \textbf{Chuyển tab} & \textbf{Background} & \textbf{Process kill} \\
\midrule
\endhead
\bottomrule
\endfoot
Video & Floating hoặc pause theo chính sách & Pause (trừ khi bật PiP) & Khôi phục video và vị trí, không tự phát \\
Music & Tiếp tục phát nếu không xung đột & Tiếp tục qua MediaSessionService & Lưu bài, vị trí và queue; khôi phục khi mở lại \\
News & Giữ vị trí cuộn & Giữ trạng thái đọc & Phục hồi article ID và vị trí đọc \\
Games & Pause + save & Pause + save & Phục hồi từ bản lưu gần nhất \\
\end{longtable}

\subsubsection{State nào được lưu ở đâu}
\label{sec:state-storage}

\begin{longtable}{>{\raggedright\arraybackslash}p{5.2cm} >{\raggedright\arraybackslash}p{9.8cm}}
\caption{Loại trạng thái và cách lưu}
\label{tab:state-storage}\\
\toprule
\textbf{Loại trạng thái} & \textbf{Cách lưu} \\
\midrule
\endfirsthead
\toprule
\textbf{Loại trạng thái} & \textbf{Cách lưu} \\
\midrule
\endhead
\bottomrule
\endfoot
Nút đang chọn, popup đang mở & Compose \texttt{rememberSaveable} \\
UI State màn hình & ViewModel + StateFlow \\
Navigation arguments & Navigation / SavedStateHandle \\
Vị trí video, bài nhạc, queue & Repository + Room/DataStore \\
Nội dung tin đã cache & Room \\
Game board, score, level & Room \\
Cài đặt mute, chất lượng media & DataStore \\
Player còn hoạt động & Playback controller / service~--~\textbf{không} lưu đối tượng Player xuống database \\
\end{longtable}

Lưu ý quan trọng: ViewModel tồn tại qua configuration changes thông thường (xoay màn hình), nhưng \textbf{không tồn tại qua process death}. Vì vậy game state và tiến độ media không thể chỉ giữ trong ViewModel~--~phải lưu xuống Room hoặc DataStore. Cũng không nên trông chờ \texttt{onDestroy()} được gọi trước khi hệ điều hành kill process; nên lưu các mốc quan trọng ngay khi trạng thái thay đổi, thay vì chỉ lưu tất cả khi app thoát.

\subsubsection{Sơ đồ vòng đời Activity}
\label{sec:activity-lifecycle}

Vòng đời Activity chuẩn của Android:

\begin{quote}
\texttt{Created} $\rightarrow$ \texttt{Started} $\rightarrow$ \texttt{Resumed} $\rightarrow$ \texttt{Paused} $\rightarrow$ \texttt{Stopped} $\rightarrow$ \texttt{Destroyed}
\end{quote}

Sự kiện chính và ứng xử:

\begin{itemize}
\item \textbf{App loses focus} (notification kéo xuống, dialog): Activity vào \texttt{Paused}~--~video pause, nhạc vẫn phát.
\item \textbf{No longer visible} (bấm Home, chuyển app): Activity vào \texttt{Stopped}~--~video và game pause + save; nhạc tiếp tục qua MediaSessionService.
\item \textbf{System reclaims memory}: process bị kill~--~khôi phục từ Room/DataStore khi mở lại.
\item \textbf{Activity finishes}: người dùng chủ động thoát hoặc hệ thống.
\end{itemize}

Các trạng thái này chỉ đại diện vòng đời \textbf{Activity}. Phát nhạc qua MediaSessionService có vòng đời riêng, không gắn cứng với MainActivity~--~khi Activity bị destroy, service vẫn chạy.

\subsubsection{Background Audio Playback}
\label{sec:background-audio}

Phân hệ Music dùng \texttt{MediaSessionService} để duy trì media session ngoài vòng đời Activity, cung cấp điều khiển từ notification, Bluetooth và màn hình khoá. Bảng~\ref{tab:audio-rules} định nghĩa hành vi mong muốn theo từng sự kiện:

\begin{longtable}{>{\raggedright\arraybackslash}p{5.4cm} >{\raggedright\arraybackslash}p{9.6cm}}
\caption{Quy tắc audio focus theo sự kiện}
\label{tab:audio-rules}\\
\toprule
\textbf{Sự kiện} & \textbf{Hành vi mong muốn} \\
\midrule
\endfirsthead
\toprule
\textbf{Sự kiện} & \textbf{Hành vi} \\
\midrule
\endhead
\bottomrule
\endfoot
Người dùng khoá màn hình & Music tiếp tục phát \\
Bấm Home & Music tiếp tục phát \\
Cuộc gọi đến & Pause theo audio focus change \\
Rút tai nghe & Pause~--~tránh phát loa ngoài bất ngờ \\
Notification Pause & Music tạm dừng \\
Notification Next & Chuyển bài tiếp theo \\
Mở Video có âm thanh & Music pause (case C1) \\
Mở game đã mute & Music tiếp tục (case C5) \\
Người dùng chủ động Stop & Dừng session nền theo chính sách \\
\end{longtable}

Khi triển khai Media3, nhóm bật xử lý Android Audio Focus qua \texttt{setAudioAttributes(..., true)}, đồng thời dùng chính sách nội bộ (bảng~\ref{tab:audio-rules}) để phối hợp các nguồn phát trong cùng ứng dụng.

\subsubsection{Ma trận chuyển phân hệ~--~đối chiếu với case C1--C9}
\label{sec:multitasking-matrix}

Bảng~\ref{tab:multitasking-matrix} là phiên bản kỹ thuật của ma trận case C1--C9 đã trình bày ở mục~\ref{sec:mvp-scope}, gắn trực tiếp với các thành phần điều phối kiến trúc:

\begin{longtable}{>{\raggedright\arraybackslash}p{3.4cm} >{\raggedright\arraybackslash}p{4.8cm} >{\raggedright\arraybackslash}p{6.8cm}}
\caption{Ma trận chuyển phân hệ: trạng thái player/game và UI}
\label{tab:multitasking-matrix}\\
\toprule
\textbf{Chuyển từ $\rightarrow$ tới} & \textbf{Player / Game State} & \textbf{UI hiển thị} \\
\midrule
\endfirsthead
\toprule
\textbf{Chuyển từ $\rightarrow$ tới} & \textbf{Player / Game} & \textbf{UI} \\
\midrule
\endhead
\bottomrule
\endfoot
Video $\rightarrow$ Tin & Video tiếp tục (case C3) & Floating Video \\
Video $\rightarrow$ Nhạc & Video pause, Nhạc play (case C1) & Mini Music Player \\
Video $\rightarrow$ Games & Video tiếp tục nếu game mute (case C4) & Floating Video + game \\
Nhạc $\rightarrow$ Tin & Nhạc tiếp tục (case C2) & Mini Music Player \\
Nhạc $\rightarrow$ Video & Nhạc pause khi video phát tiếng & Video Player \\
Nhạc $\rightarrow$ Games & Nhạc tiếp tục nếu game mute (case C5) & Game fullscreen, nhạc nền \\
Games $\rightarrow$ Tin & Game pause + save & News Screen \\
Games $\rightarrow$ Video & Game pause + save & Video Screen \\
Tin $\rightarrow$ Games & Lưu vị trí đọc & Game Screen \\
Any $\rightarrow$ Background & Pause video/game; Nhạc tiếp tục & Android background \\
\end{longtable}

Đây là \textbf{chính sách sản phẩm} nhóm chủ động lập trình, không phải hành vi Android tự động thực hiện. Mỗi chuyển tab phải được code: pause, resume, đổi player, lưu session.

\subsubsection{Quyết định kiến trúc quan trọng}
\label{sec:adr}

Bảng~\ref{tab:adr} tổng hợp các Architecture Decision Records (ADR) nhóm cần thống nhất trước khi code, để tránh các thành viên code theo cách khác nhau:

\begin{longtable}{>{\raggedright\arraybackslash}p{1.6cm} >{\raggedright\arraybackslash}p{13.4cm}}
\caption{Architecture Decision Records}
\label{tab:adr}\\
\toprule
\textbf{ADR} & \textbf{Quyết định} \\
\midrule
\endfirsthead
\toprule
\textbf{ADR} & \textbf{Quyết định} \\
\midrule
\endhead
\bottomrule
\endfoot
ADR-01 & Android Native, Kotlin + Jetpack Compose \\
ADR-02 & Android 8.0+ (API 26+), mục tiêu RAM tối thiểu 3~GB \\
ADR-03 & MVVM + 3-layer architecture + UDF (Unidirectional Data Flow) \\
ADR-04 & Feature-based modularization \\
ADR-05 & Media3 cho audio/video playback \\
ADR-06 & MediaSessionService cho Music background \\
ADR-07 & Một nguồn audio chính có tiếng tại một thời điểm; cho phép game mute SFX \\
ADR-08 & In-app Floating Player cho Video $\rightarrow$ Tin/Games \\
ADR-09 & Room cho persisted session/game state; DataStore cho preferences \\
ADR-10 & News cache và offline-first minigame \\
ADR-11 & SessionCoordinator xử lý chuyển đổi phân hệ \\
ADR-12 & Video và game mặc định pause khi app vào background; Music được tiếp tục \\
\end{longtable}

\subsubsection{Test case kiểm chứng kiến trúc}
\label{sec:arch-test-cases}

Để chứng minh kiến trúc hoạt động đúng, nhóm đưa vào báo cáo các test case thực tế sau:

\begin{longtable}{>{\raggedright\arraybackslash}p{1.4cm} >{\raggedright\arraybackslash}p{6.4cm} >{\raggedright\arraybackslash}p{7.2cm}}
\caption{Test case kiểm chứng hành vi đa nhiệm và vòng đời}
\label{tab:arch-test-cases}\\
\toprule
\textbf{Test} & \textbf{Tình huống} & \textbf{Kết quả mong đợi} \\
\midrule
\endfirsthead
\toprule
\textbf{Test} & \textbf{Tình huống} & \textbf{Kết quả} \\
\midrule
\endhead
\bottomrule
\endfoot
TC01 & Video $\rightarrow$ Tin & Video thu nhỏ, tiếp tục phát \\
TC02 & Video $\rightarrow$ Nhạc & Video pause, Nhạc bắt đầu \\
TC03 & Nhạc $\rightarrow$ Home $\rightarrow$ khoá màn hình & Nhạc tiếp tục, điều khiển từ lock screen \\
TC04 & Nhạc $\rightarrow$ Games (game mute) & Nhạc tiếp tục \\
TC05 & Game $\rightarrow$ Tin $\rightarrow$ Game & Khôi phục đúng board và score \\
TC06 & Mất mạng khi đọc Tin & Hiển thị bài đã cache \\
TC07 & Mất mạng khi stream Video & Phát buffer còn lại, sau đó trạng thái chờ/retry \\
TC08 & Có mạng trở lại & Cho phép retry hoặc resume \\
TC09 & Activity recreate (xoay màn hình) & Trạng thái cần thiết được giữ \\
TC10 & OS kill app process & Khôi phục phiên từ local storage \\
TC11 & Có cuộc gọi khi nghe Nhạc & Audio pause/duck theo chính sách \\
TC12 & Rút tai nghe & Nhạc pause \\
TC13 & Thiết bị RAM 3~GB & Không crash/OOM trong luồng demo (Quốc Huy) \\
TC14 & Bật Battery Saver & Hạn chế tác vụ không thiết yếu \\
\end{longtable}

Để kiểm chứng các tuyên bố về pin và hiệu năng, nhóm có thể dùng Android Studio Profiler, Memory Profiler và Perfetto~--~đo khi chỉ xem video so với vừa xem floating video vừa đọc tin hoặc chơi game. Các số liệu này bổ sung cho bảng chỉ số ở mục~\ref{sec:quality-pillars}.


\endgroup
